\documentclass{article}
\usepackage{/globals/de}
\usepackage{/globals/anki}
\ankiprefix{informatik}
\ankitopic{Informatik}
\ankishow
\begin{document}
Die Informatik ist die automatische Verarbeitung von Information, daher auch der Name als Kofferwort aus \textquote{Information} und \textquote{Automatik}.
\ankinote{bedeutung}{Bedeutung}{Information + Automatik}

Sie kann in vier Kernbereiche aufgeteilt werden.
\begin{description}
 \item[Theoretische Informatik] Die mathematische Theorie hinter der Informatik, der Automaten, Berechenbarkeit, Komplexität, \etc.
 \item[Angewandte Informatik] Informatik auf unterschiedliche wissenschaftliche Themen und kommerzielle Zwecke angewendet.
 \item[Praktische Informatik] Die Theorie der Algorithmen und Datenstrukturen, Funktionsweise von Programmiersprachen, Betriebssystemen, \etc.
 \item[Technische Informatik] Die Funktionsweise der Hardware der Computer.
\end{description}
\ankinote{bereiche}{Kernbereiche}{Theoretische, Angewandte, Praktische, Technische}
\ankinote{bereiche-theorie}{Kernbereich Theoretische}{Mathematische Theorie von Berechenbarkeit, Komplexität}
\ankinote{bereiche-anwendung}{Kernbereich Angewandte}{In anderen Themenbereichen}
\ankinote{bereiche-praktisch}{Kernbereich Praktische}{DSA, Sprachen, OS}
\ankinote{bereiche-technisch}{Kernbereich Technische}{Hardware}

\end{document}